% #############################################################################
% This is Chapter 4
% !TEX root = ../main.tex
% #############################################################################
% Change the Name of the Chapter i the following line
\fancychapter{Solution: ADEPT Architecture}
% The following line allows to ref this chapter
\label{chap:solution}

% NOTE: skeleton + bullet-point to-be content only (no prose yet).

% #############################################################################
\section{Threat Model}\label{sec:sol-threat}
\begin{itemize}
    \item Black-box evasion attacker (query/label access only); functional-realizability assumption --- an unrealizable attack (broken checksum, negative IAT) is out of scope by design.
    \item Attacker goal: evade the \emph{protected} downstream NIDS. ADEPT's goal: detect + purify before that NIDS ever sees the manipulated flow.
    \item Out of scope (non-goals, by design): poisoning, white-box gradient attacks, and non-protocol-compliant attacks.
\end{itemize}

% #############################################################################
\section{System Overview}\label{sec:sol-overview}
\begin{itemize}
    \item Pipeline: Packet-to-Tensor $\rightarrow$ \{TCN, RF\} $\rightarrow$ Mini-Transformer meta-learner $\rightarrow$ (if adversarial) Reconstructor $\rightarrow$ Protected NIDS.
    \item The RL agent sits outside this critical path, observing outcomes and adapting the defense over time.
    \item Two phases: Training (fit the ensemble, meta-learner, and purifier) and Deployment (the closed-loop detect/purify/adapt cycle).
\end{itemize}

% #############################################################################
\section{Dataset Construction and Packet-to-Tensor Pipeline}\label{sec:sol-pipeline}
\begin{itemize}
    \item Sources: Kitsune (Mirai, SSDP Flood), CIC-IDS2017 (DDoS, PortScan, Brute Force), CIC-DDoS-2019 (NTP/DNS/SSDP/SYN/LDAP reflection floods) $\times$ TANTGen adversarial variants (\texttt{both}/\texttt{iat}/\texttt{size}).
    \item Label convention: raw-dataset traffic = benign (upstream NIDS's job); TANTGen-generated = adversarial.
    \item nPrint bit-level per-packet encoding (1088 bits) + IPAT (z-scored) + Direction = 1090-dim token; identity-field masking (src/dst IP, ports, IP~ID).
    \item Dual emission cadence: per-packet (\texttt{flow\_packets.npy}) for the TCN; multi-scale prefix checkpoints at packet 2 then every 10 (\texttt{rf\_checkpoints.npy}, 24-dim Welford-aggregated stats) for the RF --- PortScan's $\sim$2-packet average needs an early RF checkpoint or it never gets an RF verdict.
    \item Benign-cap ratio (3:1 vs.\ adversarial) to stop the reflection-flood-dominated benign pool from swamping training; flow-atomic, source-stratified 70/15/15 split, seed~42, shared across all downstream models for fair comparison.
\end{itemize}

% #############################################################################
\section{Primary Detection: Heterogeneous Stacking Ensemble}\label{sec:sol-ensemble}
\subsection{Random Forest Base Learner}\label{subsec:sol-rf}
\begin{itemize}
    \item 24-dim aggregated prefix statistics (mean/std/min/max over IPAT, packet size, payload length/entropy, TCP flags, direction); grid-searched config (200 trees, depth~20, min\_leaf~5, balanced\_subsample).
    \item Rationale: discrete, non-differentiable boundary $\Rightarrow$ gradient masking against white-box-style perturbations even under the stated black-box threat model.
    \item Measured: ROC AUC 0.888 (per-checkpoint), Recall-adv 0.982, EDP median 2 packets, detection rate 99.69\%.
\end{itemize}
\subsection{Temporal Convolutional Network Base Learner}\label{subsec:sol-tcn}
\begin{itemize}
    \item Causal dilated TemporalBlocks $[128,128,128]$, $k=3$, receptive field 29 packets, 805K params; per-timestep $[B,T,2]$ logits --- dense, early-warning stream (verdict from packet~1).
    \item Measured: ROC AUC 0.836, EDP median 1 packet, detection rate 86.31\%.
\end{itemize}
\subsection{Attention-Based Meta-Learner (the ``Mini-Transformer'')}\label{subsec:sol-meta}
\begin{itemize}
    \item 13-dim per-packet fused token (carry-forward RF logit, dense TCN logit, rf\_age / is\_checkpoint / IPAT / direction, compact packet summary); causal mask; learnable \texttt{rf\_unknown\_logit} for pre-checkpoint NaN slots; 68K params, 2 layers / 4 heads / $d_{\text{model}}=64$.
    \item ``Beat the best base'' auxiliary loss (weight 0.5): shrinks the training confidence gap $\sim$9$\times$ but does \emph{not} move test AUROC@t (rank-bottlenecked, not loss-bottlenecked) --- reported honestly.
    \item Measured: matches/beats both base learners on the flow-level streaming axis from packet~5 onward (AUROC@5 0.834 $\rightarrow$ @max 0.845) while keeping the TCN's EDP$=1$ latency.
    \item Discussion: why RF ``wins'' on its own per-checkpoint axis but not on the comparable flow-level axis (avoid the apples-to-oranges trap).
\end{itemize}

% #############################################################################
\section{Input Purification: Reconstructor Transformer}\label{sec:sol-reconstructor}
\begin{itemize}
    \item Causal Transformer encoder + residual head ($\Delta\!\approx\!0$ default), 3 layers / 4 heads / $d_{\text{model}}=128$, 678K params; trained on 1.99M 5-tuple-matched (adversarial $\rightarrow$ clean) pairs; smooth-L1, masked over padding, length-normalized per flow.
    \item Decoder-only/causal choice over GAN/AutoEncoder justified by KV-cache streaming latency.
    \item Loss-design choice: \texttt{loss\_emphasis: none} (uniform-weight reconstruction, no attack-footprint hint) as the \textbf{main, attack-agnostic result}; the perturbation-weighted variant is an explicit off-by-default ablation.
    \item Realization step: maps purified feature deltas back onto real PCAP packets; \texttt{tshark}-validated protocol compliance (produce/validate real packets, not just tensors).
    \item Measured: MSE 0.962~$\rightarrow$~0.128 (7.5$\times$); purification gain 0.128 with \textbf{residual gap 0.000} (purified flow statistically indistinguishable from genuinely clean to the detector); 99.99\% feature-space / 99.5\% \texttt{tshark} packet compliance.
    \item Attack-dependent effect: strong on \texttt{portscan}/\texttt{ssdp\_flood}/\texttt{ddos} (\texttt{both}/\texttt{size}, gain 0.86--0.96); $\sim$0 on pure-\texttt{iat} and the four hardest reflection floods (DNS/NTP/SSDP/LDAP) --- reported as a limitation, not buried.
\end{itemize}

% #############################################################################
\section{Continuous Adaptation: Reinforcement Learning Agent}\label{sec:sol-rl}
\begin{itemize}
    \item Role in the pipeline: an RL agent sits outside the critical path, observes detection outcomes, and adapts the defense over time so it tracks an evolving attacker (high-level only).
    \item To be specified once the agent is designed: the learning formulation, what it observes, what it controls, and how the feedback loop closes --- left unspecified here on purpose.
\end{itemize}

% #############################################################################
\section{Implementation and Tools}\label{sec:sol-tools}
\begin{itemize}
    \item PyTorch (TCN / Transformer / meta-learner / reconstructor), scikit-learn (RF), Stable-Baselines3 (RL agent), Scapy + tshark/editcap/tcpdump (PCAP handling, realization, compliance checks), nPrint (bit-level encoding).
    \item Repository-layout walkthrough (config/, data/, scripts/, src/) --- brief, points to the README for detail.
\end{itemize}
